
We validated the mechanism of tri-modal reinforcement learning with dynamic feedback scheduling for code generation. Four sub-hypotheses confirmed the sequential capability building mechanism. The tri-modal framework correctly implements dynamic weight scheduling, with execution weight dominant in Phase 1 (0.$800\rightarrow 0.714)$, AI weight peaking in Phase 2 (0.545 at 50\% progress), and human weight increasing in Phase 3 (0.$400\rightarrow 0.636)$. Each phase achieved its intended objective: Phase 1 drove fastest correctness improvement ($30\times$ rate advantage), Phase 2 enabled quality refinement without correctness regression (quality +0.070, pass@1 ratio 1.032), and Phase 3 prevented execution-only collapse (conflict median 0.2468, zero samples below threshold). The 100\% gate pass rate (12/12 criteria) provides evidence the mechanism operates as predicted.

Our experiments confirm dynamic feedback scheduling is viable---the framework can be implemented, feedback collectors function, phase-specific weight patterns emerge as designed, and intended objectives are achieved. However, we explicitly do not claim performance superiority. All models used pretrained checkpoints without actual reinforcement learning training. Whether mechanism advantages translate to quantitative gains requires full-scale RL training with reward optimization. Additionally, human feedback used heuristic proxies, and we provide no comparison against optimal static weight configurations. These limitations are transparent and acceptable for proof-of-concept validation.

The conceptual contribution extends beyond code generation. The question shifts from ''which feedback signal?'' to ''which signal when?'' This opens research directions: curriculum learning over feedback modality, learned weight schedules through meta-learning, adaptive schedules responding to training dynamics, and extension to additional feedback modalities.

Immediate next steps include full RL training validation with performance comparison, static versus dynamic ablation, real human annotation collection, and cross-domain application. We have shown the mechanism is viable and testable---future work must determine whether it yields performance advantages at scale.
